\chapter{How Computers Execute Programs}

\section{Purpose}
Establish the CPU and memory execution model that motivates later compilation details.

\section{Outline}
\begin{enumerate}
  \item CPU execution model and architectural state
  \item Instruction cycle: fetch, decode, execute
  \item Registers and register file roles
  \item Memory model and addressing
  \item Stack and function call mechanics
  \item Calling conventions overview
  \item Transition to ABI discussion
\end{enumerate}

\section{Content Plan}
\begin{itemize}
  \item CPU execution model: instructions update architectural state and memory.
  \item Instruction cycle: control flow, sequencing, and side effects.
  \item Registers: general purpose, PC, SP, flags, callee vs caller saved.
  \item Memory model: byte addressing, alignment, endianness.
  \item Stack mechanics: frame layout, prologue and epilogue, return address.
  \item Calling conventions: argument passing, return values, stack cleanup.
  \item ABI bridge: why a stable contract is required.
\end{itemize}

\section{Key Terms}
\begin{itemize}
  \item instruction pointer
  \item stack frame
  \item calling convention
  \item ABI
\end{itemize}

\section{Diagrams Planned}
\begin{figure}[h]
  \centering
  \fbox{\parbox{0.9\linewidth}{CPU state and instruction cycle diagram placeholder.}}
  \caption{Planned CPU state and instruction cycle diagram.}
\end{figure}

\begin{figure}[h]
  \centering
  \fbox{\parbox{0.9\linewidth}{Stack frame layout diagram placeholder.}}
  \caption{Planned stack frame layout diagram.}
\end{figure}

\section{Experiment}
See how-computers-execute-programs-exp/how-computers-execute-programs-exp.tex in this chapter.
